\documentclass[12pt]{extarticle}

\usepackage{tabularx}
\usepackage{amsmath}
\usepackage{extsizes}
\usepackage{graphicx}
\usepackage{amssymb}
\usepackage[english,russian]{babel}
\usepackage[utf8]{inputenc}
\usepackage[a4paper]{geometry}
\geometry{top=20mm,bottom=30mm,left=20mm,right=15mm}
\usepackage[T1,T2A]{fontenc}
\usepackage{setspace}
\usepackage{indentfirst}
\usepackage{caption}
\renewcommand{\baselinestretch}{1.5}
\captionsetup[figure]{name={Рис.},labelsep=endash}

\begin{document}

\begin{center}
  \textbf{МИНОБРНАУКИ РОССИИ} \\
  \textbf{Санкт-Петербургский государственный} \\
  \textbf{электротехнический университет} \\
  \textbf{«ЛЭТИ» им. В.И. Ульянова (Ленина)} \\
  \textbf{Кафедра САПР} \\
\end{center}

\vspace{0.25\textheight}

\begin{center}
\textbf{ОТЧЕТ} \\
\textbf{по лабораторной работе №3 (9)} \\
\textbf{по дисциплине «Автобусы»} \\
\textbf{Тема: КРИТИЧЕСКИЕ ИСПЫТАНИЕ И КРАШТЕСТ} \\
\end{center}

\vspace*{\fill}

\begin{center}
\begin{tabularx}{\textwidth}{ l X r }
 Студент гр. 5352 & \rule{8cm}{0.2mm} & Иванов И.И. \\
 Преподаватель & \rule{8cm}{0.2mm} & Иванов И.И. \\
\end{tabularx}
\end{center}

\vspace{3em}

\begin{center}
Санкт-Петербург\\
\the\year
\end{center}

\thispagestyle{empty}

\pagebreak

АВТОБУС

\end{document}
